\section{Fase 1: Definición del proyecto}
\label{sec:definicion-proyecto}

El punto de partida del proyecto fue una reunión con el tutor en abril de 2025, en la que le comenté posibles ideas de trabajo.
En esa primera toma de contacto surgieron un par de ideas más, las cuales finalmente se acabaron descartando en nuestra segunda reunión,
ya que decidimos adaptarnos un poco más a la situacion actual con el tema de las paginas webs y automatizaciones.


La primera idea consistía en una aplicación web que, a partir de datos de APIs hospitalarias, ayudara a los pacientes a encontrar el centro 
médico más adaptado a sus necesidades según una serie de parámetros. 
La segunda idea ampliaba esa primera hacia un asistente de mudanza. Imaginese que alguien quiere instalarse en una ciudad, la aplicación
recogería todos los datos de las APIs publicas y un parseo de los datos proporcionados por agencias y otras empresas, de tal forma que 
introduciendo el usuario sus requisitos, se le organizaría un posible plan de vida en la zona.
Ambas ideas tenían en común la lectura y procesamiento de APIs publicas, pero el alcance era demasiado limitado, lo que limitaba el valor del 
proyecto. 


En la segunda reunion, realizada en Julio de 2025, llegamos a una idea más acotada, una aplicacion web Django encargada de procesar los datos
de cualquier API publica y generar automaticamente un sitio web personalizado. Mas tarde pensé en posibles nombres para el proyecto, hasta que 
al final dimos con el nombre de WebBuilder.


En ese momento aún estaba liado con los examenes, pero en los ratos libres iba recogiendo informacion y creando las posibles fases del sistema:
validación de la URL, procesamiento de los datos recogidos, preferencias sobre el diseño, etc. David valoró positivamente estos pasos, lo que 
me dio mas energía para seguir adelante con el proyecto.


Durante los meses de verano el proyecto comenzo poco a poco creando las partes mas sencillas para ir dandole forma: la pagina de inicio, el 
sistema de login y registro de usuarios. Aunque en el futuro todo esto cambiaría, la base mas  o menos se mantendría.

% Poco a poco conforme el proyecto avanzo fui teniendo problemas para poder hacer de verdad  algo personalizado y diferente, era muy dificil procesar los datos de todas las APIs y que no saliesen datos genericos, ahi le surgió a David una idea, la cual dio el giro completo del proyecto y lo convirtio en algo innovador.

% Ibamos a incorporar un modulo de Inteligencia Artificial, la idea no solo era procesarle los datos y que el eligiese los mas adecuados, sino construir un sistema en el que la IA fuese el motor principal y actuase como un desarrollador guiado con prompting engenieering. Pero este cambio estara detallado mas adelante.

Durante los primeros meses siguientes no hubo un gran desarrollo de codigo, sino que dedicamos este periodo a definir con mas detalle el flujo
de la aplicacion, estuvimos investigando tecnologias y diseñando procesos que deberia seguir el sistema. Este trabajo previo fue muy importante
para el correcto desarrollo guiado de la aplicacion.


